\documentclass[french]{article}

\usepackage[utf8]{inputenc}
\usepackage[T1]{fontenc}
\usepackage{lmodern}
\usepackage{geometry}
\usepackage{graphicx}
\usepackage{babel}
\usepackage{amsmath}
\usepackage{amsthm}
\usepackage{amssymb}
\usepackage{hyperref}
\usepackage{stmaryrd}
\usepackage{enumitem}
\usepackage[most]{tcolorbox}
\usepackage{dsfont}
\usepackage{multirow}
\usepackage{etoc}
\usepackage{titlesec}
\usepackage{esint}
\usepackage{cancel}
\usepackage{mathdots}
\usepackage{indentfirst}
\usepackage{lastpage}
\usepackage{fancyhdr}

\pagestyle{fancy}
\fancyhf{}
\renewcommand{\headrulewidth}{0pt}
\renewcommand{\footrulewidth}{0pt}
\fancyhead[L,R]{}
\fancyfoot[R]{\thepage \ /\ \pageref{LastPage}}

\fancypagestyle{plain}{\fancyhead{}\renewcommand{\headrule}{}}

\setlength{\parindent}{0pt}

\setlist[itemize]{label=$\bullet$}
\setlist{before=\leavevmode, leftmargin=*}

\renewcommand{\P}{\mathbb{P}}
\newcommand{\N}{\mathbb{N}}
\newcommand{\Z}{\mathbb{Z}}
\newcommand{\Q}{\mathbb{Q}}
\newcommand{\R}{\mathbb{R}}
\newcommand{\C}{\mathbb{C}}
\newcommand{\K}{\mathbb{K}}
\renewcommand{\L}{\mathbb{L}}
\newcommand{\U}{\mathbb{U}}
\newcommand{\st}{^*}
\newcommand{\tq}{\middle|}
\newcommand{\inv}{^{-1}}
\newcommand{\E}{\mathbb{E}}
\newcommand{\F}{\mathbb{F}}
\newcommand{\V}{\mathbb{V}}
\renewcommand{\ker}{\text{Ker}}
\newcommand{\im}{\text{Im}}
\newcommand{\card}{\text{Card}}
\newcommand{\Tr}{\text{Tr}}
\newcommand{\Vect}{\text{Vect}}
\newcommand{\rg}{\text{rg}}
\newcommand{\vertiii}[1]{{\left\vert\kern-0.25ex\left\vert\kern-0.25ex\left\vert #1 \right\vert\kern-0.25ex\right\vert\kern-0.25ex\right\vert}}
\renewcommand{\o}{\text{o}}
\renewcommand{\O}{\text{O}}
\renewcommand{\d}{\text{d}}
\newcommand{\Id}{\text{Id}}


\theoremstyle{plain}
\newtheorem*{ind}{Indication}

\theoremstyle{definition}

\newtheorem*{ex}{Exemple}
\newtheorem*{rmq}{Remarque}
\newtheorem*{notation}{Notation}
\newtheorem*{rappel}{Rappel}
\newtheorem*{terminologie}{Terminologie}
\newtheorem*{attention}{Attention}
\newtheorem*{conséquence}{Conséquence}

\theoremstyle{remark}
\newtheorem*{app}{Application}

\newtcolorbox[auto counter]{dfn}[1][]{
	enhanced,
	colback=white,
	colframe=black,
	sharp corners,
	boxrule=0.8pt,
	fonttitle=\bfseries,
	colbacktitle=white,
	coltitle=black,
	attach boxed title to top left={yshift=-1mm,xshift=-0.5mm},
	boxed title style={size=minimal, right=2pt, sharp corners, colframe=white},
	title={Définition \thetcbcounter \ifstrempty{#1}{}{ (#1)}},
	before skip=0.5em,
	after skip=0.5em
}

\newtcolorbox[auto counter]{exo}[1][]{
	enhanced,
	arc=1mm,
	colback=white,
	colframe=black,
	coltitle=black,
	fonttitle=\bfseries,
	boxrule=0.8pt,
	shadow={1.2pt}{-1.2pt}{0pt}{black!50},
	attach title to upper={\ },
	title={Exercice \thetcbcounter\ifblank{#1}{}{ #1}},
	before skip=0.5em,
	after skip=0.5em,
	left=2mm, right=2mm, top=1mm, bottom=1mm
}

\newtcolorbox{met}[1][]{
	enhanced,
	arc=1mm,
	colback=white,
	colframe=black,
	coltitle=black,
	fonttitle=\bfseries,
	boxrule=0.8pt,
	shadow={1.2pt}{-1.2pt}{0pt}{black!50},
	attach title to upper={\ },
	title={Point méthode \ifblank{#1}{}{#1}},
	before skip=1em,
	after skip=1em,
	left=2mm, right=2mm, top=1mm, bottom=1mm
}

\newcounter{thmcount}

\newtcolorbox[use counter=thmcount]{thm}[1][]{
	enhanced,
	colback=white,
	colframe=black,
	sharp corners,
	left skip=0mm,
	boxrule=1.3pt,
	fonttitle=\bfseries,
	colbacktitle=white,
	coltitle=black,
	attach boxed title to top left={yshift=-1mm,xshift=-0.5mm},
	boxed title style={size=minimal, right=2pt, sharp corners, colframe=white},
	title={Théorème \thethmcount \ifstrempty{#1}{}{ (#1)}},
	before skip=0.5em,
	after skip=0.5em
}

\newtcolorbox[use counter=thmcount]{prop}[1][]{
	enhanced,
	colback=white,
	colframe=black,
	sharp corners,
	boxrule=0.8pt,
	fonttitle=\bfseries,
	colbacktitle=white,
	coltitle=black,
	attach boxed title to top left={yshift=-1mm,xshift=-0.5mm},
	boxed title style={size=minimal, right=2pt, sharp corners, colframe=white},
	title={Proposition \thethmcount \ifstrempty{#1}{}{ (#1)}},
	before skip=0.5em,
	after skip=0.5em
}

\newtcolorbox[use counter=thmcount]{cor}[1][]{
	enhanced,
	colback=white,
	colframe=black,
	sharp corners,
	boxrule=0.8pt,
	fonttitle=\bfseries,
	colbacktitle=white,
	coltitle=black,
	attach boxed title to top left={yshift=-1mm,xshift=-0.5mm},
	boxed title style={size=minimal, right=2pt, sharp corners, colframe=white},
	title={Corollaire \thethmcount \ifstrempty{#1}{}{ (#1)}},
	before skip=0.5em,
	after skip=0.5em
}

\newtcolorbox[use counter=thmcount]{lem}[1][]{
	enhanced,
	colback=white,
	colframe=black,
	sharp corners,
	boxrule=0.5pt,
	fonttitle=\bfseries,
	colbacktitle=white,
	coltitle=black,
	attach boxed title to top left={yshift=-1mm,xshift=-0.5mm},
	boxed title style={size=minimal, right=2pt, sharp corners, colframe=white},
	title={Lemme \thethmcount \ifstrempty{#1}{}{ (#1)}},
	before skip=0.5em,
	after skip=0.5em
}

\addto\captionsfrench{\renewcommand*{\contentsname}{Intégration sur un intervalle quelconque}}

\title{Intégration sur un intervalle quelconque}
\date{}

\begin{document}
\tableofcontents	

\newpage
\maketitle

\section{Intégrale généralisée}

Dans ce chapitre, on étend la notion d'intégrale, déjà connue pour les fonctions continues par morceaux sur un segment, à certaines fonctions continues par morceaux sur un intervalle $I$ quelconque de $\R$ et à valeurs dans $\K=\R$ ou $\K=\C$.\\
Les intervalles de $\R$ considérés dans ce chapitre sont tous d'intérieur non vide.

\subsection{Fonction continue par morceaux}

Rappelons la définition d'une fonction continue par morceaux sur un segment et étendons-la à un intervalle quelconque.

\begin{dfn}
	\begin{itemize}
		\item Une fonction $f:[a,b]\to\K$ est dite \textbf{continue par morceaux sur le segment $[a,b]$} de $\R$ s'il existe une subdivision $(a_0,\ldots,a_n)$ de ce segment telle que, pour tout $i\in\llbracket0,n-1\rrbracket$, la restriction de $f$ à $]a_i,a_{i+1}[$ possède un prolongement continu sur le segment $[a_i,a_{i+1}]$?
		\item Une fonction $f:I\to\K$ est dite \textbf{continue par morceaux sur l'intervalle $I$} si sa restriction à tout segment inclus dans $I$ est continue par morceaux.
	\end{itemize}
\end{dfn}

\begin{notation}
	On note $\mathcal{CM}(I,\K)$ l'ensemble des fonctions continues par morceaux sur l'intervalle $I$ à valeurs dans $\K$.
\end{notation}

\begin{rmq}
	L'ensemble $\mathcal{CM}(I,\K)$ contient les fonctions constantes sur $I$, est stable par combinaison linéaire et produit : c'est une sous-algèbre de $\K^I$. On a immédiatement l'inclusion $\mathcal{C}(I,\K)\subset\mathcal{CM}(I,\K)$.
\end{rmq}

\begin{ex}
	\begin{enumerate}
		\item La fonction partie entière est une fonction continue par morceaux sur $\R$.
		\item La fonction inverse est continue par morceaux sur $]0,+\infty[$ mais ne possède pas de prolongement continue par morceaux sur $[0,+\infty[$ puisqu'elle n'est pas bornée au voisinage de $0$.
	\end{enumerate}
\end{ex}

\begin{attention}
	Il n'y a pas de théorème concernant la composition ni le quotient de fonctions continues par morceaux !
\end{attention}

\begin{rappel}
	Lorsque $f$ est une fonction continue sur $I$, pour $a\in I$, la fonction : $$\begin{array}{lrcl}
	F:&I&\longrightarrow&\K\\
	&x&\longmapsto&\displaystyle\int_{a}^{x}f(t)\d t
\end{array}$$ est dérivable, de dérivée $f$. On en déduit, en particulier, que si $f$ est positive, alors $F$ est croissante.
\end{rappel}
Voyons ce que ces résultats deviennent pour des fonctions continues par morceaux.

\begin{prop}%prop
	Soit $f$ une fonction continue par morceaux et positive sur $I$. Alors la fonction :
	$$\begin{array}{lrcl}
		F:&I&\longrightarrow&\R\\
		&x&\longmapsto&\displaystyle\int_{a}^{x}f(t)\d t
	\end{array}$$
	est croissante sur $I$.
\end{prop}

\begin{prop}%prop
	Soit $f:I\to\K$ une fonction continue par morceaux et $a\in I$. La fonction :
	$$\begin{array}{lrcl}
		F:&I&\longrightarrow&\K\\
		&x&\longmapsto&\displaystyle\int_{a}^{x}f(t)\d t
	\end{array}$$
	est continue sur $I$. Elle est :
	\begin{itemize}
		\item dérivable à droite en tout point de $x$ différent de $\sup(I)$, avec $F_d'(x)=f(x^+)$ ;
		\item dérivable à gauche en tout point de $x$ différent de $\inf(I)$, avec $F_g'(x)=f(x^-)$.
	\end{itemize}
	En particulier, $F$ est dérivable en tout point de continuité de $f$ avec $F'(x)=f(x)$.
\end{prop}

\begin{rmq}
	Par extension, on dira que $F:x\mapsto\displaystyle\int_{a}^{x}f(t)\d t$ est une \textbf{primitive} de $f$.
\end{rmq}

\subsection{Définitions}

\subsubsection*{Cas où $I=[a,b[$}

Ici, on considère $a\in\R$ et $b\in\overline{\R}$ tels que $a<b\leqslant+\infty$.

\begin{dfn}
	Soit $f\in\mathcal{CM}([a,b[,\K)$.
	On dit que l'\textbf{intégrale $\displaystyle\int_{[a,b[}f$ converge} si la fonction $x\mapsto\displaystyle\int_{[a,x]}f$ admet une limite finie en $b$. Dans ce cas, on note $\displaystyle\int_{[a,b[}f$ ou $\displaystyle\int_{[a,b[}f(t)\d t$ cette limite.\\
	Dans le cas contraire, on dit que l'intégrale $\displaystyle\int_{[a,b[}f$ \textbf{diverge} (et on ne la note pas).
\end{dfn}

\begin{terminologie}
	\begin{itemize}
	\item Une telle intégrale est appelée \textbf{intégrale généralisée}. On parle aussi d'\textbf{intégrale impropre}.
	\item Pour une intégrale généralisée, son caractère convergent ou divergent est appelé sa nature.
\end{itemize}
\end{terminologie}

\begin{prop}[Caractère local de la nature d'une intégrale]%prop
	Soit $f\in\mathcal{CM}([a,b[,\K)$ ainsi que $c\in[a,b[$.\\
	Les intégrales $\displaystyle\int_{[a,b[}f$ et $\displaystyle\int_{[c,b[}f$ sont de même nature et, si elles convergent, on a :
	$$\int_{[a,b[}f=\int_{[a,c]}f+\int_{[c,b[}f.$$
\end{prop}
\begin{proof}
	On utilise la relation de Chasles et la définition de la convergence d'une intégrale.
\end{proof}

\begin{cor}%cor
	Soit $f\in\mathcal{CM}([a,b[,\K)$. Si l'intégrale converge, $\displaystyle\int_{[a,b[}f$ converge, alors $\displaystyle\int_{[x,b[}f\xrightarrow[x\to b]{}0$.
\end{cor}

\begin{rmq}
	On peut noter une analogie entre l'étude des intégrales généralisées sur un intervalle $[a,+\infty[$ et celle des séries numériques.
\end{rmq}

\begin{terminologie}
	Par analogie avec les séries, l'intégrale $\displaystyle\int_{[a,x]}f$ est appelée \textbf{intégrale partielle} et, en cas de convergence, l'intégrale $\displaystyle\int_{[x,b[}f$ est appelée \textbf{reste} de l'intégrale $\displaystyle\int_{[a,b[}f$.
\end{terminologie}


Le résultat suivant assure que, pour une fonction $f$ continue par morceaux sur le segment $[a,b]$, il n'y a pas de différence entre son intégrale sur le segment $[a,b]$ (au sens de la première année) et son intégrale sur l'intervalle semi-ouvert $[a,b[$ (au sens que nous venons de définir).

\begin{prop}%prop
	Soit $f\in\mathcal{CM}([a,b],\K)$. Alors l'intégrale généralisée $\displaystyle\int_{[a,b[}f$ converge et l'on a :
	$$\int_{[a,b[}f=\int_{[a,b]}f$$
\end{prop}

\begin{notation}
	Dans la suite, l'intégrale $\displaystyle\int_{[a,b[}f$ est notée $\displaystyle\int_{a}^{b}f$. Dans le cas où $f$ est continue par morceaux sur le segment $[a,b]$, le résultat précédent assure la cohérence de cette notation avec celle déjà utilisée en première année.
\end{notation}

\begin{prop}%prop
	Soit $f\in\mathcal{C}([a,b[,\K)$ et $F$ une primitive de $F$. L'intégrale $\displaystyle\int_{a}^{b}f$ converge si, et seulement si, $F$ possède une limite finie en $b$. Dans ce cas, on a : 
	$$\int_{a}^{b}f=\underset{b}{\lim}F-F(a)$$ et la fonction définie sur $[a,b[$ par $x\mapsto\displaystyle\int_{x}^{b}f$ est dérivable, de dérivée $-f$.
\end{prop}

\begin{notation}
	Si $\underset{b}{\lim}F$ existe, on note $[F]_a^b$ l'expression $\underset{b}{\lim}F-F(a)$.
\end{notation}

\begin{ex}
	\begin{enumerate}
		\item L'intégrale $\displaystyle\int_{0}^{+\infty}e^{-t}\d t$ converge et vaut $\displaystyle\int_{0}^{+\infty}e^{-t}\d t=1$.
		\item Montrons que l'intégrale $\displaystyle\int_{1}^{+\infty}\frac{\d t}{t(1+t^2)}$ converge et calculons sa valeur. %par cdv x=t^2, $0,5\ln(x^2/(1+x^2)) en est une primitive, on trouve finalement $\ln(2)/2$.
		\item L'intégrale $\displaystyle\int_{0}^{+\infty}\cos(t)\d t$ diverge. %sin n'admet pas de limite
	\end{enumerate}
\end{ex}

\begin{attention}
	A l'instar des séries numériques, la condition $f\xrightarrow[+\infty]{}0$ n'est pas suffisante pour que l'intégrale $\displaystyle\int_{a}^{+\infty}f$ converge, comme le montre l'exemple suivant.
\end{attention}

\begin{ex}
	Considérons $f:t\mapsto\displaystyle\frac{1}{t}$ sur $[1,+\infty[$. $f\xrightarrow[+\infty]{}0$, mais l'intégrale $\displaystyle\int_{1}^{+\infty}f$ diverge car :
	$$\forall x\geqslant1,\ \int_{1}^{x}\frac{\d t}{t}=\ln(x)\xrightarrow[x\to +\infty]{}+\infty.$$
\end{ex}

\begin{attention}
	En revanche, contrairement au cas des séries numériques, la convergence de l'intégrale $\displaystyle\int_{a}^{+\infty}f$ n'entraîne le fait que $f$ tende vers $0$ en $+\infty$, ni même que $f$ est bornée au voisinage de $+\infty$, comme on le verra dans un exemple dans la suite du chapitre.
\end{attention}

\begin{met}
	Pour $N\in\N$ et $a>0$, on a d'après la relation de Chasles : $$\sum_{n=1}^{N}\int_{na}^{(n+1)a}f=\int_{a}^{(N+1)a}f$$ Par conséquent, si l'intégrale $\displaystyle\int_{a}^{+\infty} f$ converge alors la série $\displaystyle\sum \int_{na}^{(n+1)a}f$ converge. Ainsi, on peut utilise la divergence d'une série pour prouver la divergence d'une intégrale. Attention, la réciproque est fausse.
\end{met}

\begin{ex}
	Montrons que l'intégrale $\displaystyle\int_{\pi}^{+\infty}\left\lvert\frac{\sin(t)}{t}\right\rvert \d t$ diverge.
\end{ex}

\begin{attention}
	La réciproque est fausse : la convergence de la série $\displaystyle\sum \int_{na}^{(n+1)a}f$ n'implique pas la convergence de l'intégrale $\displaystyle\int_{a}^{+\infty}f$, comme le montre l'exemple suivant.
\end{attention}

\begin{ex}
	L'intégrale $\int_{0}^{+\infty}\cos(t)\d t$ diverge mais $\displaystyle\sum \int_{n\pi}^{2(n+1)\pi}\cos(t)\d t$ est convergente, car son terme général est nul.
\end{ex}

Voir néanmoins le point méthode de la section $1.3$.

\subsubsection*{Cas où $I=]a,b]$}

Ici, on considère $a\in\overline{\R}$ et $b\in\R$ tels que $-\infty\leqslant a<b$.

\begin{dfn}
	Soit $f\in\mathcal{CM}(]a,b],\K)$.
	On dit que l'\textbf{intégrale $\displaystyle\int_{]a,b]}f$ converge} si la fonction $x\mapsto\displaystyle\int_{[x,b]}f$ admet une limite finie en $a$. Dans ce cas, on note $\displaystyle\int_{]a,b]}f$ ou $\displaystyle\int_{]a,b]}f(t)\d t$ cette limite.\\
	Dans le cas contraire, on dit que l'intégrale $\displaystyle\int_{]a,b]}f$ \textbf{diverge} (et on ne la note pas).
\end{dfn}

\begin{ex}
	$\displaystyle\int_{]0,1]}\ln(t)\d t=-1$.
\end{ex}

\begin{exo}
	Déterminer la nature et la valeur éventuelle de l'intégrale $\displaystyle\int_{0}^{\frac{\pi}{2}}\frac{\cos(u)}{\sqrt{\sin(u)}}\d u$.
\end{exo}

\begin{rmq}
	On adapte sans difficulté au cas d'un intervalle semi-ouvert $]a,b]$ les propositions énoncées précédemment dans le cas d'un intervalle du type $[a,b[$.
\end{rmq}

\begin{notation}
	L'intégrale $\displaystyle\int_{]a,b]}f$ est aussi notée $\displaystyle\int_{a}^{b}f$.
\end{notation}

\subsubsection*{Cas d'un intervalle ouvert $]a,b[$}

Dans cette section, on considère $(a,b)\in\overline{\R}^2$ tel que $-\infty\leqslant a\leqslant b\leqslant+\infty$.

\begin{lem}%lem
	Soit $f\in\mathcal{CM}(]a,b[,\K)$.\\
	S'il existe $c\in]a,b[$ tel que les deux intégrales $\displaystyle\int_{]a,c]}f$ et $\displaystyle\int_{[c,b[}f$ convergent, alors, pour tout réel $c'\in]a,b[$, les intégrales $\displaystyle\int_{]a,c']}f$ et $\displaystyle\int_{[c',b[}f$ convergent aussi, et :
	$$\int_{]a,c]}f+\int_{[c,b[}f=\int_{]a,c']}f+\int_{[c',b[}f.$$
\end{lem}
\begin{proof}
	Utiliser le caractère local de la nature d'une intégrale.
\end{proof}

On peut donc poser la définition suivante.

\begin{dfn}
	Soit $f\in\mathcal{CM}(]a,b[,\K)$. On dit que l'intégrale $\displaystyle\int_{]a,b[}f$ \textbf{converge} s'il existe un réel $c\in]a,b[$ tels que les deux intégrales $\displaystyle\int_{]a,c]}f$ et $\displaystyle\int_{[c,b[}f$ convergent. On pose alors : 
	$$\int_{]a,b[}f=\int_{]a,c]}f+\int_{[c,b[}f.$$ Dans le cas contraire, on dit que l'intégrale $\displaystyle\int_{]a,b[}f$ \textbf{diverge}.
\end{dfn}

\paragraph{Compatibilité avec les définitions déjà données}
\begin{itemize}
	\item Si $f\in\mathcal{CM}(]a,b],\K)$, il découle de la définition précédente et du caractère local de la nature d'une intégrale que les intégrales $\displaystyle\int_{]a,b[}f$ et $\displaystyle\int_{]a,b]}f$ ont même nature et, en cas de convergence, sont égales. On a le même résultat si $f\in\mathcal{CM}([a,b[,\K)$.
	\item De même, un autre résultat précédent s'étend : si $f\in\mathcal{CM}([a,b],\K)$, l'intégrale $\displaystyle\int_{]a,b[}f$ converge et on a $\displaystyle\int_{]a,b[}f=\int_{[a,b}f$.
\end{itemize}

\begin{notation}
	La discussion précédente justifie alors la notation $\displaystyle\int_{a}^{b}f$ utilisée pour désigner cette intégrale.
\end{notation}

\begin{ex}
	Si $\displaystyle\int_{-\infty}^{+\infty}f$ converge, alors $\displaystyle\int_{-x}^{x}f\xrightarrow[x\to+\infty]{}\int_{-\infty}^{+\infty}f$.
\end{ex}

\begin{attention}
	La réciproque du résultat obtenue par l'exemple précédent est fausse : comme le montre l'exemple suivant, l'existence de $\displaystyle\underset{x\to+\infty}{\lim}\int_{-x}^{x}f(t)\d t$ n'entraîne pas nécessairement la convergence de l'intégrale $\displaystyle\int_{-\infty}^{+\infty}f$.
\end{attention}

\begin{ex}
	Pour tout $x\in\R$, on a $\displaystyle\int_{-x}^{x}t\ \d t=0$, donc $\displaystyle\int_{-x}^{x}t\ \d t\xrightarrow[x\to+\infty]{}0$, mais l'intégrale $\displaystyle\int_{-\infty}^{+\infty}t\ \d t$ n'est pas convergente puisque $\displaystyle\int_{0}^{x}t\ \d t\xrightarrow[x\to+\infty]{}+\infty$.
\end{ex}

\begin{met}
	Si $F$ est une primitive de $f$, alors l'intégrale $\displaystyle\int_{a}^{b}f$ converge si, et seulement si, $F$ admet des limites finies en $a$ et en $b$, et l'on a alors : $$\int_{a}^{b}f=\lim_{b}F-\lim_{a}F$$ quantité alors notée $[F(t)]_a^b$.
\end{met}

\begin{ex}
	$\displaystyle\int_{-\infty}^{+\infty}\frac{\d t}{1+t^2}$.
\end{ex}

\subsection{Cas des fonctions réelles à valeurs positives}

De la même façon que pour les séries numériques, on étudie le cas des fonctions positives. La notion d'intégrabilité permettra ensuite de se ramener à ce cas.\\
On traite dans ce paragraphe le cas des intégrales généralisées sur l'intervalle semi-ouvert $[a,b[$, le cas $]a,b]$ étant analogue et le cas $]a,b[$ se ramenant aux deux précédents.

\begin{prop}%prop
	Soit $a\in\R$ et $b\in\overline{\R}$ tels que $a<b\leqslant+\infty$, ainsi que $f\in\mathcal{CM}([a,b[,\R)$ à valeurs \textit{positives}.
	\begin{enumerate}
		\item L'intégrale $\displaystyle\int_{a}^{b}f$ converge si, et seulement si, $x\mapsto\displaystyle\int_{a}^{x}f$ est majorée sur $[a,b[$.
		\item L'intégrale $\displaystyle\int_{a}^{b}f$ diverge si, et seulement si, $\displaystyle\int_{a}^{x}f\xrightarrow[x\to b]{}+\infty$. On pose alors : $\displaystyle\int_{a}^{b}f=+\infty$.
	\end{enumerate}
\end{prop}
\begin{proof}
	On utilise la croissance de la fonction $x\mapsto\displaystyle\int_{a}^{x}f$ et le théorème de la limite monotone.
\end{proof}

\begin{rmq}
	\begin{itemize}
		\item Les conclusions sur la nature de l'intégrale subsistent si $f$ n'est positive qu'au voisinage de $b$, car $x\mapsto\displaystyle\int_{a}^{x}f$ est alors croissante au voisinage de $b$.
		\item L'adaptation de la proposition précédente au cas d'une fonction $f\in\mathcal{CM}(]a,b],\R)$ positive est la suivante : l'intégrale $\displaystyle\int_{a}^{b}f$ converge si, et seulement si, la fonction $x\mapsto\displaystyle\int_{x}^{b}f$ est majorée, et si ce n'est pas le cas, alors $\displaystyle\int_{x}^{b}f\xrightarrow[x\to a]{}+\infty$.
	\end{itemize}
\end{rmq}

\begin{prop}[Comparaison d'intégrales de fonctions positives]%prop
	Soit $a\in\R$ et $b\in\overline{\R}$ tels que $a<b\leqslant+\infty$, ainsi que $(f,g)\in\mathcal{CM}([a,b[,\R)^2$ à valeurs \textit{positives}.
	Si $f\leqslant g$ sur $[a,b[$, alors on a l'inégalité suivant dans $\overline{\R}_+$ : $$\int_{a}^{b}f\leqslant\int_{a}^{b}g.$$
	En particulier, si l'intégrale $\displaystyle\int_{a}^{b}g$ converge, alors l'intégrale $\displaystyle\int_{a}^{b}f$ converge.
\end{prop}
\begin{proof}
	On utilise la caractérisation de la convergence de l'intégrale d'une fonction positive (proposition précédente) pour écrire : $$\int_{a}^{x}f\leqslant\int_{a}^{x}g\leqslant\int_{a}^{b}g$$ et on passe à la limite sur $x$.
\end{proof}

\begin{met}
	Soit $a>0$ et $f\in\mathcal{CM}([a,+\infty[,\R_+)$. Puisqu'on a pour tout $N\in\N$ : $$\sum_{n=1}^{N}\int_{na}^{(n+1)a}f=\int_{a}^{(N+1)a}f$$ on obtient dans $\overline{\R_+}$, par passage à la limite et positivité : $$\sum_{n=1}^{+\infty}\int_{na}^{(n+1)a}f=\int_{a}^{+\infty}f$$ En particulier, la série $\displaystyle\sum \int_{na}^{(n+1)a}f$ et l'intégrale l'intégrale $\displaystyle\int_{a}^{+\infty} f$ ont même nature.
\end{met}

\begin{ex}
	Fonction continue en dents-de-scie dont les pics sont de largeur $1/n$ et de hauteur $n$.
\end{ex}

\begin{exo}
	Pour tout $n\in\N\st$, on pose $u_n=\displaystyle\sum_{k=1}^{n}\frac{1}{k}-\ln(n)$. On définit $f$ sur $]0,+\infty[$ par $f(t)=\displaystyle\frac{t-\lfloor t\rfloor}{t^2}$.
	\begin{enumerate}
		\item Montrer que $\displaystyle\int_{1}^{+\infty}f(t)\d t$ est bien définie dans $[0,+\infty]$ et : $$\int_{1}^{+\infty}f(t)\d t=\sum_{j=1}^{+\infty}\left(\ln(j+1)-\ln(j)-\frac{1}{j+1}\right)<+\infty.$$
		\item Montrer que la suite $(u_n)_{n\in\N\st}$ converge vers $1-\displaystyle\int_{1}^{+\infty}f(t)\d t$.
	\end{enumerate}
\end{exo}

\subsection{Intégrales de référence}

Comme pour les séries numériques, nous disposons d'un certain nombre d'exemples de référence permettant, grâce au théorème de comparaison, de déterminer la nature de beaucoup d'autres intégrales.

\begin{prop}%prop
	Soit $\alpha\in\R$. L'intégrale $\displaystyle\int_{0}^{+\infty}e^{-\alpha t}\ \d t$ converge si, et seulement si, $\alpha>0$.
\end{prop}

\begin{prop}[Intégrales de Riemann]%prop
	Soit $\alpha\in\R$.
	\begin{itemize}
		\item L'intégrale $\displaystyle\int_{1}^{+\infty}\frac{\d t}{t^\alpha}$ converge si, et seulement si, $\alpha>1$.
		\item L'intégrale $\displaystyle\int_{0}^{1}\frac{\d t}{t^\alpha}$ converge si, et seulement si, $\alpha<1$.
		\item Plus généralement, étant donné $(a,b)\in\R^2$ tel que $a<b$, les intégrales $\displaystyle\int_{a}^{b}\frac{\d t}{(b-t)^\alpha}$ et $\displaystyle\int_{a}^{b}\frac{\d t}{(t-a)^\alpha}$ convergent si, et seulement si, $\alpha<1$.
	\end{itemize}
\end{prop}

\begin{ex}
	\begin{enumerate}
		\item Les intégrales $\displaystyle\int_{1}^{+\infty}\frac{\d t}{t^2}$, $\displaystyle\int_{0}^{1}\frac{\d t}{\sqrt{t}}$ et $\displaystyle\int_{0}^{1}\frac{\d t}{\sqrt{1-t}}$ convergent.
		\item Les intégrales $\displaystyle\int_{1}^{+\infty}\frac{\d t}{t}$, $\displaystyle\int_{0}^{1}\frac{\d t}{t}$ et $\displaystyle\int_{1}^{3}\frac{\d t}{(3-t)^2}$ divergent.
	\end{enumerate}
\end{ex}

\section{Propriétés de l'intégrale}

Les propositions suivantes étendent à un intervalle $I$ quelconque des propriétés déjà connues sur un segment. Rappelons que $I$ désigner un intervalle de $\R$ d'\textit{intérieur non vide}.

\paragraph{Convention} Par extension, lorsque $I$ est un segment et que $f$ est continue par morceaux sur $I$, on dit encore que l'intégrale $\displaystyle\int_If$ converge.

\subsection{Linéarité}

\begin{prop}[Linéarité]%prop
	L'ensemble des fonctions $f\in\mathcal{CM}(I,\K)$ telles que $\displaystyle\int_I f$ converge est un sous-espace vectoriel de $\mathcal{CM}(I,\K)$ sur lequel l'application $f\mapsto\displaystyle\int_I f$ est une forme linéaire.
\end{prop}
\begin{proof}
	Cela résulte de la définition de la convergence d'une intégrale, de la linéarité de l'intégrale sur un segment et des théorèmes d'opérations sur les limites.
\end{proof}

\begin{prop}[Fonctions à valeurs complexes]%prop
	Étant donnée $f\in\mathcal{CM}(I,\C)$, l'intégrale $\displaystyle\int_I f$ converge si, et seulement si, les deux intégrales $\displaystyle\int_I\text{Re}(f)$ et $\displaystyle\int_I\text{Im}(f)$ convergent, et l'on a alors : $$\int_{I}f=\int_I\text{Re}(f)+i\int_I\text{Im}(f).$$
\end{prop}
\begin{proof}
	Cela résulte de la définition de la convergence d'une intégrale et des propriétés des limites des fonctions à valeurs complexes.
\end{proof}

\subsection{Positivité, croissance}

On considère ici des fonctions à valeurs réelles.

\begin{prop}[Positivité]%prop
	Soit $f\in\mathcal{CM}(I,\R)$ telle que $\displaystyle\int_I f$ converge. Si $f$ est positive sur $I$, alors $\displaystyle\int_If\geqslant0$.
\end{prop}

\begin{prop}[Croissance]%prop
	Soit $(f,g)\in\mathcal{CM}(I,\R)^2$ tel que $\displaystyle\int_I f$ et $\displaystyle\int_I g$ convergent. Si $f\leqslant g$ sur $I$, alors $\displaystyle\int_If\leqslant\int_I g$.
\end{prop}
\begin{proof}
	On applique la positivité de l'intégrale à la fonction $g-f$ en utilisant la linéarité de l'intégrale.
\end{proof}

\begin{prop}[Fonction continue positive et d'intégrale nulle]
	Soit $f\in\mathcal{C}(I,\R)$ continue à valeurs positives. Si $\displaystyle\int_If=0$, alors $f=0$.
\end{prop}
\begin{proof}
	Dans le cas où $I$ est de la forme $[a,b[$, on dérive $x\mapsto\displaystyle\int_{a}^{x}f(t)\ \d t$.
\end{proof}

\begin{met}
	Si $f$ est une fonction continue positive et non nulle sur $I$, alors $\displaystyle\int_I f>0$.
\end{met}

\begin{rmq}
	Hors minoration ou calcul explicite, c'est la seule méthode dont on dispose pour montrer qu'une intégrale est strictement positive.
\end{rmq}

\subsection{Relation de Chasles}

\begin{prop}%prop
	Soit $f\in\mathcal{CM}(I,\K)$ telle que $\displaystyle\int_If$ converge. Si $J$ est un sous-intervalle d'intérieur non vide de $I$, alors $\displaystyle\int_Jf$ converge.
\end{prop}
\begin{proof}
	Le cas $J=I$ étant trivial, supposons $J\subsetneq I$.
	Si aucune des deux extrémités $\alpha$ et $\beta$ de $J$ n'est une extrémité ouvert de $I$, la fonction $f$ est définie et continue par morceaux sur le segment $[\alpha,\beta]$, et il a déjà été indiqué qu'alors l'intégrale $\displaystyle\int_Jf$ converge. Supposons donc qu'au moins l'un des extrémités ouvertes de $J$ soit une extrémité ouverte de $I$.\\
	\textbf{Cas $I=[a,b[$.} Dans ce cas, on a $J=[c,b[$ avec $a\leqslant c<b$, et la convergence de $\displaystyle\int_Jf$ résulte du caractère local de la convergence d'une intégrale.\\
	\textbf{Cas $I=]a,b]$.} Se traite de manière analogue au précédent.\\
	\textbf{Cas $I=]a,b[$} Dans ce cas, $J=]a,c]$ ou $J=[c,b[$ avec $a<c<b$ ; la convergence de l'intégrale $\displaystyle\int_Jf$ provient alors de la définition de la convergence de l'intégrale $f$ sur l'intervalle ouvert $]a,b[$.
\end{proof}

\begin{conséquence}
	Supposons que $\displaystyle\int_I f$ converge. Si $x<y$ sont des points ou des extrémités de $I$, alors on a $]x,y[\subset I$, donc l'intégrale $\displaystyle\in_{]x,y[}f$, notée aussi $\displaystyle\int_{x}^{y}f$, converge.
\end{conséquence}

\begin{notation}
	Plus généralement, si $\displaystyle\int_I f$ converge, la notation $\displaystyle\int_{x}^{y}f$, déjà définie dans le cas où $x<uy$, est étendue au cas général ainsi :
$$\int_{x}^{y}f=-\int_{y}^{x}f\quad\text{si}\quad x>y\quad\quad\text{et}\quad\quad\int_{x}^{y}f=0\quad\text{si}\quad x=y.$$
\end{notation}

\begin{prop}[Relation de Chasles]%prop
	Soit $f\in\mathcal{CM}(I,\K)$ telle que $\displaystyle\int_If$ converge. Si $x$, $y$ et $z$ sont trois points ou extrémités de $I$, alors les intégrales $\displaystyle\int_{x}^{z}f$, $\displaystyle\int_{x}^{y}f$ et $\displaystyle\int_{y}^{z}f$ convergent, et l'on a : 
	$$\int_{x}^{z}f=\int_{x}^{y}f+\int_{y}^{z}f.$$
\end{prop}

\subsection{Intégrabilité}

Dans cette section, $I$ désigne un intervalle de $\R$ d'intérieur non vide.

\begin{dfn}
	On dit que $f\in\mathcal{CM}(I,\K)$ est \textbf{intégrable} sur $I$, ou que l'intégrale $\displaystyle\int_I f$ \textbf{converge absolument}, si l'intégrale $\displaystyle\int_I \lvert f\rvert$ converge.
\end{dfn}

\begin{notation}
	On note $L^1(I,\K)$ l'ensemble des fonctions intégrables sur $I$.
\end{notation}

\begin{prop}%prop
	$L^1(I,\K)$ est un sous-espace vectoriel de $\mathcal{CM}(I,\K)$.
\end{prop}
\begin{proof}
	Utiliser l'inégalité triangulaire et le théorème de comparaison pour les fonctions à valeurs positives.
\end{proof}

\begin{rmq}
	\begin{itemize}
		\item Si $f\in\mathcal{CM}(I,\K)$ est une fonction \textit{positive}, alors l'intégrabilité de $f$ équivaut à la convergence de l'intégrale $\displaystyle\int_I f$.
		\item En particulier, les fonctions des exemples de référence étant toutes positives, leur convergence équivaut à leur intégrabilité.
	\end{itemize}
\end{rmq}

\begin{thm}%thm
	Si $f\in L^1(I,\K)$, alors $\displaystyle\int_If$ converge.
\end{thm}
\begin{proof}
	On étudier d'abord le cas des fonctions $f$ à valeurs réelles, en utilisant les fonctions $f^+=\displaystyle\frac{\lvert f\rvert+f}{2}$ et $f^-=\displaystyle\frac{\lvert f\rvert-f}{2}$, puis on en déduit le cas des fonctions à valeurs complexes en écrivant $f=\text{Re}(f)+i\text{Im}(f)$.
\end{proof}

\begin{thm}[Inégalité triangulaire]%thm
	Pour tout $f\in L^1(I,\K)$, on a $\displaystyle\left\lvert\int_If\right\rvert\leqslant\int_I \lvert f\rvert$.
\end{thm}

\begin{prop}[Caractère local de l'intégrabilité]%prop
	Sot $a\in\R$ et $b\in\overline{\R}$ tels que $a<b\leqslant+\infty$. Soit $f\in\mathcal{CM}([a,b[,\K)$. Alors, pour tout $c\in[a,b[$, la fonction $f$ est intégrable sur $[a,b[$ si, et seulement si, elle l'est sur $[c,b[$.
\end{prop}
\begin{proof}
	C'est une conséquence immédiate du caractère local de la nature d'une intégrale.
\end{proof}

\begin{conséquence}
	L'intégrabilité de $f$ sur $[a,b[$ ne dépend que du comportement local de $f$ au voisinage de $b$.
\end{conséquence}

\begin{terminologie}
	Pour $f\in\mathcal{CM}([a,b[,\K)$, on dira que $f$ est \textbf{intégrable en $b$} si $f$ est intégrable sur $[a,b[$.
\end{terminologie}

\subsection{Intégrabilité et comparaison}

\begin{thm}[Théorème de comparaison]%thm
	Soit $(a,b)\in\R\times\overline{\R}$ tel que $a<b\leqslant+\infty$ et $(f,g)\in\mathcal{CM}([a,b[,\K)^2$.
	\begin{itemize}
		\item Si $f\underset{b}{=}\O(g)$ et si $g$ est intégrable sur $[a,b[$, alors $f$ est intégrable sur $[a,b[$.
		\item Si $f\underset{b}{\sim}g$, alors l'intégrabilité de $f$ sur $[a,b[$ équivaut à celle de $g$.
	\end{itemize}
\end{thm}
\begin{proof}
	Cela résulte de la définition de l'intégrabilité, du théorème de comparaison pour les fonctions à valeurs positives et du caractère local de l'intégrabilité.
\end{proof}

\begin{attention}
	Le théorème de comparaison énoncé ici concerne l'\textit{intégrabilité} des des fonctions, c'est-à-dire la convergence \textit{absolue} de leurs intégrales. On peut très bien avoir $f\underset{b}{\sim}g$ sans que les intégrales $\displaystyle\int_{a}^{b}f$ et $\displaystyle\int_{a}^{b}g$ soient de même nature. Se référer à un exemple et une remarque de la section sur la semi-convergence.
\end{attention}

\begin{rmq}
	Ce qui précède (en particulier le théorème de comparaison) a été énoncé dans le cas où $I=[a,b[$ mais s'adapte naturellement au cas $I=]a,b]$.
\end{rmq}

\begin{exo}[Intégrales de Bertrand]
	On étudie ici, en fonction de $(\alpha,\beta)\in\R^2$ la nature de l'\textbf{intégrale de Bertrand} : $$B_{\alpha,\beta}=\int_{2}^{+\infty}\frac{\d t}{t^\alpha(\ln t)^\beta}.$$
	\begin{enumerate}
		\item Établir la convergence de $B_{\alpha,\beta}$ pour $\alpha>1$.
		\item Établir la divergence de $B_{\alpha,\beta}$ pour $\alpha<1$.
		\item Faire l'étude de la nature de $B_{\alpha,\beta}$ pour $\alpha=1$.
	\end{enumerate}
\end{exo}

\begin{met}
	Pour étudier l'intégrabilité d'une fonction continue par morceaux sur un intervalle $I$, on peut s'intéresser au comportement asymptotique de $f$ au voisinage des "bornes ouvertes" de $I$, puis utiliser le théorème de comparaison. En particulier, si l'intervalle $I$ est ouvert, on étudiera séparément les deux bornes.
\end{met}

\begin{ex}
	\begin{enumerate}
		\item Soit $x>0$. Montrons la convergence de l'intégrale $\displaystyle\int_{0}^{+\infty}t^{x-1}e^{-t}\d t$ (\textbf{fonction Gamma} d'Euler). La fonction $f:t\mapsto t^{x-1}e^{-t}$ est continue sur $]0,+\infty[$.\\
		\textbf{Étude en $0$.} Quand $t$ tend vers $0$, on a $f(t)\sim\displaystyle\frac{1}{t^{1-x}}$. Comme $1-x<1$ la fonction $t\mapsto\displaystyle\frac{1}{t^{1-x}}$ est intégrable (intégrale de Riemann), donc par comparaison, $f$ est intégrable en $0$.\\
		\textbf{Étude en $+\infty$.} Par croissances comparées, on a : $$t^{x-1}e^{-t}\underset{+\infty}{=}\O\left(\frac{1}{t^2}\right).$$
		Comme $2>1$, la fonction $t\mapsto\displaystyle\frac{1}{t^2}$ est intégrable en $+\infty$, donc $f$ aussi par comparaison. \\
		En conclusion, la fonction $f$ est intégrable sur $]0,+\infty[$ donc son intégrale converge.
		
		\item Déterminons la nature de l'intégrale $\displaystyle\int_{0}^{+\infty}\sin(t)\ln\left(\frac{t^2+2}{t^2+1}\right)\d t$.\\
		La fonction $f:t\mapsto\displaystyle\sin(t)\ln\left(\frac{t^2+2}{t^2+1}\right)$ est continue sur $[0,+\infty[$. On a par calcul asymptotique : $$f(t)\underset{+\infty}{=}\O\left(\frac{1}{t^2}\right).$$
		Comme $2>1$, la fonction $t\mapsto\displaystyle\frac{1}{t^2}$ est intégrable en $+\infty$, donc $f$ aussi par théorème de comparaison. Par conséquent, $f$ est intégrable sur $[0,+\infty[$ donc son intégrale converge.
		
		\item Déterminons la nature de l'intégrale $\displaystyle\int_{0}^{+\infty}\frac{\cos(x)}{\sqrt{x}+x^2}\d x$.\\
		La fonction $f:x\mapsto\displaystyle\frac{\cos(x)}{\sqrt{x}+x^2}$ est continue sur $]0,+\infty[$.\\
		\textbf{Étude en $0$.} On a $f(x)\underset{x\to 0}{\sim}\frac{1}{\sqrt{x}}$. Comme $\displaystyle\frac12<1$, la fonction $\displaystyle x\mapsto\frac{1}{\sqrt{x}}$ est intégrable en $0$ donc $f$ aussi par théorème de comparaison.\\
		\textbf{Étude en $+\infty$.} On a $f(x)\displaystyle\underset{x\to+\infty}{=}\O\left(\frac{1}{x^2}\right)$. Comme $2>1$, la fonction $x\mapsto\displaystyle\frac{1}{x^2}$ est intégrable en $+\infty$, donc $f$ aussi par théorème de comparaison.\\
		En conclusion, la fonction $f$ est intégrable sur $]0,+\infty[$ donc son intégrale converge.
	\end{enumerate}
\end{ex}

\begin{met}
	On a les différents résultats.
	\begin{itemize}
		\item Si une fonction $f$ est bornée au voisinage d'un point $b\in\R$, alors on a $f\underset{b}{=}\O(1)$. La fonction constante égale à $1$ étant intégrable sur tout intervalle borné, on en déduit que $f$ est intégrable en $b$.
		\item En particulier, si $f\in\mathcal{CM}([a,b[,\K)$ admet une limite finie en $b\in\R$, alors $f$ est intégrable.
	\end{itemize}
\end{met}

\begin{exo}
	Dans ce dernier cas, en prolongeant $f$ par continuité en $b$ et une fonction $\tilde{f}$, a-t-on $\displaystyle\int_{[a,b[}f=\int_{[a,b]}\tilde{f}$ ? %peu de chances que ça ait du sens, mais si cela en a, oui
\end{exo}

\begin{attention}
	Le caractère \textit{fini} de $b$ est essentiel, comme le montre l'exemple de la fonction $t\mapsto\displaystyle\frac{1}{t}$, bornée et non intégrable sur $[1,+\infty[$.
\end{attention}

\begin{ex}
	La fonction $f:t\mapsto\sin\left(\frac{1}{t^2}\right)$ est intégrable sur $]0,+\infty[$ car bornée au voisinage de $0$ et équivalente à $t^{-2}$ positive et intégrable en $+\infty$.
\end{ex}

\begin{ex}
	Soit $x\leqslant0$. Montrons la divergence de l'intégrale $\displaystyle\int_{0}^{+\infty}t^{x-1}e^{-t}\d t$.\\
	La fonction $t\mapsto t^{x-1}e^{-t}$ est continue par morceaux sur $]0,+\infty[$. Comme elle est de plus positive, la convergence de son intégrale équivaut à son intégrabilité. L'équivalent : $$t^{x-1}e^{-t}\underset{t\to0}{\sim}t^{x-1},$$ prouve, par comparaison avec un exemple de Riemann divergent, que la fonction $t\mapsto t^{x-1}e^{-t}$ n'est pas intégrable en $0$.\\
	Conclusion : l'intégrale $\displaystyle\int_{0}^{+\infty}t^{x-1}e^{-t}\ \d t$.
\end{ex}

\subsubsection*{Utilisation d'une série}

\begin{ex}
	\begin{enumerate}
		\item Pour $\alpha\in\R_+\st$, on pose $\begin{array}{lrcl}
			f_\alpha:&[\pi,+\infty[&\longrightarrow&\R\\
			&t&\longmapsto&\displaystyle\frac{\sin(t)}{t^{\alpha}}
		\end{array}$.
		\begin{itemize}
			\item Il est clair, par comparaison, que $f_\alpha$ est intégrable sur $[\pi,+\infty[$ si $\alpha>1$.
			\item Pour $\alpha\leqslant1$, on peut utiliser la série de terme général $_n=\displaystyle\int_{n\pi}^{(n+1)\pi}f_\alpha(t)\ \d t$.
		\end{itemize}
		\item Convergence de l'intégrale $\displaystyle\int_{0}^{+\infty}\frac{\d t}{1+t^{\alpha}\sin^2(t)}$.
	\end{enumerate}
\end{ex}

\section{Calcul d'intégrales}

\subsection{Intégration par parties}

\paragraph{Convention} Si $h$ est une fonction continue sur un intervalle d'extrémités $a$ et $b$, avec $-\infty\leqslant a<b\leqslant+\infty$, on dit que le \textbf{crochet $[h]_a^b$ converge} si $h$ possède des limites \textit{finies} en $a$ et en $b$, et l'on pose alors $[h]_a^n=\underset{b}{\lim}h-\underset{a}{\lim}h$.

\begin{rmq}
	\begin{itemize}
		\item Dans le cas d'un segment $[a,b]$, par continuité de $h$, on a $\underset{a}{\lim}h=h(a)$ et $\underset{b}{\lim}h=h(b)$, et donc simplement $[h]_a^b=h(b)-h(a)$.
		\item Dans le cas d'un intervalle semi-ouvert, par exemple $[a,b[$, alors le crochet $[h]_a^b$ converge si, et seulement si, $h$ possède une limite finie en $b$, et l'on a alors $[h]_a^b=\underset{b}{\lim}h-h(a)$.
	\end{itemize}
\end{rmq}

\begin{prop}[Théorème d'intégration par parties]%prop
	Soit $f$ et $g$ deux fonctions de classe $\mathcal{C}^1$ sur un intervalle d'extrémités $a$ et $b$, avec $-\infty\leqslant a<b\leqslant+\infty$.\\
	Si le crochet $[fg]_a^b$ converge, alors les deux intégrales $\displaystyle\int_{a}^{b}f'g$ et $\displaystyle\int_{a}^{b}fg'$ sont de même nature et, en cas de convergence, on a : $$\int_{a}^{b}f'g=[fg]_a^b-\int_{a}^{b}fg'.$$
\end{prop}

\begin{ex}
	Calcul de l'intégrale $\displaystyle\int_{1}^{+\infty}\frac{\arctan(t)}{t^2}\d t$.
\end{ex}

\begin{met}
	En pratique, on peut commencer par intégrer par parties sur les primitives, puis passer à la limite. Il y a deux buts possibles :
	\begin{itemize}
		\item justifier une convergence ;
		\item effectuer un calcul.
	\end{itemize}
\end{met}

\begin{ex}
	Étudions la convergence de $I=\displaystyle\int_{0}^{1}\frac{\ln(1-t^2)}{t^2}\d t$ et calculons son éventuelle valeur. %prolonger par continuité en 0, dériver le $\ln$, on trouve $-2\ln(2)$
\end{ex}

\begin{exo}
	\begin{enumerate}
		\item Montrer que pour tout $n\in\N\st$, l'intégrale $I_n=\displaystyle\int_{0}^{+\infty}\frac{\d t}{(1+t^2)^n}$.
		\item A l'aide d'une intégration par parties, exprimer $I_{n+1}$ en fonction de $I_n$.
		\item Calculer $I_n$.
	\end{enumerate}
\end{exo}

\subsection{Changement de variables}

Soit $(a,b,\alpha,\beta)\in\overline{\R}^4$ tel que $-\infty\leqslant a<b\leqslant+\infty$ et $-\infty\leqslant\alpha<\beta\leqslant+\infty$.

\begin{prop}[Théorème de changement de variable (cas croissant)]%prop
	Soit $f\in\mathcal{CM}(]a,b[,\K)$ et $\varphi:]\alpha,\beta[\to]a,b[$ une fonction de classe $\mathcal{C}^1$, strictement croissante et bijective. Alors, les intégrales $\displaystyle\int_{a}^{b}f(t)\d t$ et $\displaystyle\int_{\alpha}^{\beta}f(\varphi(u))\varphi'(u)\d u$ sont de même nature et, en cas de convergence :
	$$\int_{a}^{b}f(t)\d t=\int_{\alpha}^{\beta}f(\varphi(u))\varphi'(u)\d u.$$
\end{prop}
\begin{proof}
	Remarquer qu'en notant $F$ une primitive de $f$, alors $F\circ\varphi$ est une primitive de $\varphi'\times(f\circ\varphi)$, et utiliser un point méthode précédent.
\end{proof}

\begin{rmq}
	Dans le cas d'un intervalle $[a,b[$, on pourra se ramener à la situation de la proposition précédente, puisque si $\displaystyle\int_{[a,b[}f$ converge, il en est de même pour $\displaystyle_{]a,b[}$ et les deux intégrales sont égales. Même principe pour un intervalle $]a,b]$.
\end{rmq}

\begin{ex}
	Montrons l'égalité $\displaystyle\int_{0}^{+\infty}e^{-t^2}\d t=\frac{1}{2}\int_{0}^{+\infty}\frac{e^{-u}}{\sqrt{u}}\d u$.
\end{ex}

\begin{prop}[Théorème de changement de variable (cas décroissant)]%prop
	Soit $f\in\mathcal{CM}(]a,b[,\K)$ et $\varphi:]\alpha,\beta[\to]a,b[$ une fonction de classe $\mathcal{C}^1$, strictement décroissante et bijective. Alors, les intégrales $\displaystyle\int_{a}^{b}f(t)\d t$ et $\displaystyle\int_{\alpha}^{\beta}f(\varphi(u))\varphi'(u)\d u$ sont de même nature et, en cas de convergence :
	$$\int_{a}^{b}f(t)\d t=-\int_{\alpha}^{\beta}f(\varphi(u))\varphi'(u)\d u.$$
\end{prop}
\begin{proof}
	Le principe est le même que pour la proposition précédente, en notant bien que cette fois on a : $$\underset{\alpha}{\lim}\varphi=b\quad\text{et}\quad\underset{\beta}{\lim}\varphi=a.$$
\end{proof}

\begin{rmq}
	\begin{itemize}
		\item Dans chacun des deux cas, on a 	$\displaystyle\int_{a}^{b}f(t)\d t=\int_{\alpha}^{\beta}f(\varphi(u))\lvert\varphi'(u)\rvert \d u.$
		\item Comme $\varphi'$ est de signe constant, en appliquant le théorème de changement de variable à $\lvert f\rvert$, on obtient que l'intégrabilité de $f$ sur $]a,b[$ sur $]a,b[$ équivaut à celle de $(f\circ\varphi)\times\varphi'$ sur $]\alpha,\beta[$.
	\end{itemize}
\end{rmq}

\begin{ex}
	Montrons la convergence de l'intégrale $I=\displaystyle\int_{0}^{+\infty}\frac{\ln(t)}{(1+t)^2}\d t$ et calculons sa valeur.
\end{ex}

\subsection{Semi-convergence}

Il se peut qu'une fonction $f\in\mathcal{CM}(I,\C)$ ne soit pas intégrable, mais que l'intégrale $\displaystyle\int_I f$ converge. on dit alors que cette dernière intégrale est \textbf{semi-convergente}.


Comme pour les séries semi-convergentes, l'étude de la convergence d'une intégrale non absolument convergente ne peut pas se faire complètement à l'aide des théorèmes de comparaison, puisque ces derniers ne concernent que des fonctions intégrables.


Il va donc falloir transformer l'intégrale pour pouvoir étudier sa convergence. Un changement de variable ne sert à rien en général puisque, d'après la remarque précédente, on ne pourra pas obtenir une intégrale absolument convergente en partant d'une intégrale semi-convergente par changement de variable.


L'outil principal est l'\textit{intégration par parties}, avec en général comme objectif de transformer l'intégrale en faisant apparaître une intégrale absolument convergente.

\begin{ex}
	\begin{enumerate}
		\item Soit $\alpha\in]0,1]$. Étudions la convergence de l'intégrale $\displaystyle\int_{1}^{+\infty}\frac{e^{it}}{t^{\alpha}}$.
		\item Convergence des deux intégrales $\displaystyle\int_{1}^{+\infty}\frac{\cos(t)}{t^{\alpha}}\d t$ et $\displaystyle\int_{1}^{+\infty}\frac{\sin(t)}{t^{\alpha}}\d t$.
	\end{enumerate}
\end{ex}

\begin{rmq}
	Pour l'étude de certaines intégrales non absolument convergentes, on pourra, comme pour les séries, effectuer un développement asymptotique, le dernier terme écrit étant :
	\begin{itemize}
		\item ou bien une fonction intégrable ;
		\item ou bien de signe constant (dans le cas où la fonction est à valeurs réelles).
	\end{itemize}
\end{rmq}

\begin{ex}
	Déterminons la nature des intégrales $\displaystyle\int_{0}^{+\infty}\ln\left(1+\frac{\sin(t)}{\sqrt{t}}\right)\d t$ et $\displaystyle\int_{0}^{+\infty}\frac{\sin(t)}{\sqrt{t}}\d t$.
\end{ex}

\begin{rmq}
	Remarquons que nous avons ainsi obtenu deux fonctions continues par morceaux sur $[1,+\infty[$, équivalentes en $+\infty$, mais dont les intégrales ne sont pas de même nature.
\end{rmq}

\begin{exo}
	Montrer la convergence de l'intégrale $\displaystyle\int_{1}^{+\infty}\sin\left(\frac{\sin(t)}{\sqrt{t}}\right)\d t$.
\end{exo}
	
\section{Intégration des relations de comparaison}

Ici, les fonctions sont définies sur un intervalle $[a,b[$ avec $-\infty<a<b\leqslant+\infty$, et l'on fait une étude locale au voisinage de $b$.

On procéderait de la même façon dans le cas d'un intervalle $]a,b]$, pour une étude au voisinage de $a$.

\subsection{Cas convergent : comparaison des restes}

\begin{prop}[Intégration des relations de comparaison (cas convergent)]%prop
	Soit $f\in\mathcal{CM}([a,b[,\K)$, ainsi que $\varphi\in\mathcal{CM}([a,b[,\R])$ \textit{positive} et \textit{intégrable}.
	\begin{itemize}
		\item Si $f\underset{b}{=}\O(\varphi)$, alors $f$ est intégrable sur $[a,b[$ et $\displaystyle\int_{x}^{b}f\underset{x\to b}{=}\O\left(\int_{x}^{b}\varphi\right)$.
		\item Si $f\underset{b}{=}\o(\varphi)$, alors $f$ est intégrable sur $[a,b[$ et $\displaystyle\int_{x}^{b}f\underset{x\to b}{=}\o\left(\int_{x}^{b}\varphi\right)$.
		\item Si $f\underset{b}{\sim}\varphi$, alors $f$ est intégrable sur $[a,b[$ et $\displaystyle\int_{x}^{b}f\underset{x\to b}{\sim}\int_{x}^{b}\varphi$.
	\end{itemize}
\end{prop}

\begin{rmq}
	L'énoncé précédent, donnant un résultat au voisinage de $b$, reste valable si la fonction $\varphi$ n'est positive qu'au voisinage de $b$.
\end{rmq}

\begin{ex}
	\begin{enumerate}
		\item Existence, équivalent puis développement asymptotique de $\displaystyle\int_{x}^{+\infty}\frac{\arctan(t)}{t\sqrt{t}}\d t$ lorsque $x$ tend vers $+\infty$.
		\item Équivalent simple de $\arccos$ en $1$.
	\end{enumerate}
\end{ex}

\begin{exo}
	Pour tout $n\in\N\st$, on pose $u_n=\displaystyle\int_{n}^{+\infty}e^{-t^2}\d t$.
	\begin{enumerate}
		\item Montrer, pour tout $n\in\N\st$, la convergence de l'intégrale définissant $u_n$.
		\item Déterminer un équivalent simple de la suite $(u_n)_{n\in\N\st}$. On pourra effectuer une intégration par parties.
	\end{enumerate}
\end{exo}

\begin{met}
	On peut appliquer le théorème si la fonction à laquelle on compare est de signe constant et intégrable.
\end{met}

\subsection{Cas divergent : comparaison des intégrales partielles}

\begin{prop}[Intégration des relations de comparaison (cas divergent)]%prop
	Soit $f\in\mathcal{CM}([a,b[,\K)$, ainsi que $\varphi\in\mathcal{CM}([a,b[,\R])$ \textit{positive} et \textit{non intégrable}.
	\begin{itemize}
		\item Si $f\underset{b}{=}\O(\varphi)$, alors $f$ est intégrable sur $[a,b[$ et $\displaystyle\int_{a}^{x}f\underset{x\to b}{=}\O\left(\int_{a}^{x}\varphi\right)$.
		\item Si $f\underset{b}{=}\o(\varphi)$, alors $f$ est intégrable sur $[a,b[$ et $\displaystyle\int_{a}^{x}f\underset{x\to b}{=}\o\left(\int_{a}^{x}\varphi\right)$.
		\item Si $f\underset{b}{\sim}\varphi$, alors $f$ est intégrable sur $[a,b[$ et $\displaystyle\int_{a}^{x}f\underset{x\to b}{\sim}\int_{a}^{x}\varphi$.
	\end{itemize}
\end{prop}

\begin{ex}
	\begin{enumerate}
		\item Soit $f\in\mathcal{CM}(\R_+,\K)$ et $l\in\K$. On suppose que $f(x)\xrightarrow[x\to+\infty]{}l$.\\
		Alors on a $f(x)-l\underset{x\to+\infty}{=}o(1)$ et la fonction constante égale à $1$ est positive et non intégrable sur $[0,+\infty[$. Par intégration des relations de comparaison (cas divergent), on en déduit : $$\int_{0}^{x}(f(t)-l)\d t\underset{x\to+\infty}{=}\o\left(\int_{0}^{x}1\ \d t\right)\quad\text{donc}\quad\int_{0}^{x}f(t)\ \d t\underset{x\to+\infty}{=}lx+\o(x).$$
		\item Déterminons, à l'aide d'une intégration par parties, un équivalent simple de $\displaystyle Li(x)=\int_{2}^{x}\frac{\d t}{\ln(t)}$ quand $x$ tend vers $+\infty$.
	\end{enumerate}
\end{ex}

\begin{exo}
	Pour tout $x>0$, on pose $f(x)=\displaystyle\int_{0}^{+\infty}\frac{e^{-t}}{x+t}\d t$.
	\begin{enumerate}
		\item Montrer que pour tout $x>0$, l'intégrale définissant $f(x)$ est convergente.
		\item Déterminer un équivalent simple de $f(x)$ quand $x$ tend vers $0$.
	\end{enumerate}
\end{exo}

\begin{met}
	On peut appliquer le théorème si la fonction à laquelle on compare est de signe constant et non intégrable.
\end{met}

\end{document}